% language=us

\environment context-2024-style

\logo [XEON]    {XEON}
\logo [OPENSUSE]{OpenSUSE}
\logo [QEMU]    {QEMU}

\startdocument
  [title={\TeX},
   banner={into the current landscape},
   location={context\enspace {\bf 2024}\enspace meeting}]


\starttitle[title=Over 30 years of \TEX]

The 1995 4all\TEX\ \DVD\ shipped with \CONTEXT\ and (last week) installing it on
the old \IBM\ 380 Laptop was still possible. A simple \quotation {hello world} document
processed surprisingly fast (a few seconds). We're talking of 223 pages with
standard layout and fonts.

This is huge Em\TeX, some \DVI\ bitmap output turned \POSTSCRIPT, viewed with
\GHOSTSCRIPT. Around that time we used \DVIPSONE\ and outline \POSTSCRIPT\ fonts
already with \DVIWINDO\ and \ACROBAT\ on \MSDOS. Replicating that nowadays is non
trivial. A 386 \CPU, 640 KB memory and a 512 MB disk.

\stoptitle

\starttitle[title=Some 10 years later]

A 2005 version of \CONTEXT\ on an old Dell Inspiron 8000 laptop processes 1000
\quote {tufte in a macro} samples in 7.5 seconds. We're talking \PDFTEX,
\TEXEXEC\ and \TEXUTIL.

\starttyping
\def\Tufte{We live ...}

\starttext \dorecurse{1000}{\Tufte\par} \stoptext
\stoptyping

\stoptitle

\starttitle[title=Again about 20 years later]

The same test on my Dell 7520 laptop (Xeon 1505M v6) takes:

\starttabulate[|l|l|l|l|]
\NC 1.2 seconds \NC \PDFTEX     \NC \MKII \NC           \NC \NR
\NC 5.7 seconds \NC \XETEX      \NC \MKII \NC           \NC \NR
\NC 3.3 seconds \NC \LUATEX     \NC \MKIV \NC           \NC \NR
\NC 2.9 seconds \NC \LUAMETATEX \NC \MKXL \NC aka \LMTX \NC \NR
\stoptabulate

On a slightly newer Dell 7540 laptop (i7 9750H) it takes 2.2 seconds with
an August 2024 \LMTX.

\stoptitle

\starttitle[title=Today 2024]

Feel free to benchmark:

\starttyping
\starttext
    \dorecurse{1000}{\samplefile{tufte}\par}
\stoptext
\stoptyping

Run it a few times in order to populate the OS file cache. It runs the same
speed on a similar \LINUX\ laptop (also a 7520 with Xeon).

It can't be worse than on a \QEMU\ \RISCV\ emulator which takes 46.1 seconds on that
same 7520 laptop (15 times slower).

\stoptitle

\starttitle[title=Some things stay]

\startitemize
    \startitem A fifteen year old car is just as comfortable as a modern one. \stopitem
    \startitem Next years mobile phone is not substantially better than last years. \stopitem
    \startitem My more than 5 year old laptop feels quite okay in 2024. \stopitem
    \startitem Books evolved into perfection and are still valued. \stopitem
  % \startitem Bad things also stay, take dictators, warlords and injustice. \stopitem
    \startitem Decent \EPUB\ devices will eventually arrive, real digital paper. \stopitem
\stopitemize

\stoptitle

\starttitle[title=Some things change purpose]

\startitemize
    \startitem Wood and brick windmills are monuments but still occasionally used. \stopitem
    \startitem Church towers still are landmarks. \stopitem
    \startitem Water towers are obsolete although on top of buildings \unknown\ (USA). \stopitem
    \startitem Castles are obsolete but still impress tourists and serve history. \stopitem
\stopitemize

\stoptitle

\starttitle[title=Some things disappear]

\startitemize
    \startitem Huge companies seldom age well. \stopitem
    \startitem Semaphores populated the landscape for a while but disappeared fast. \stopitem
    \startitem Programming languages become obsolete. \stopitem
    \startitem Programs disappear, technology is ditched. \stopitem
    \startitem Ice cellars and stores are solidly replaced by refrigerators. \stopitem
    \startitem Fossil fuel cars will be replaced by battery powered ones. \stopitem
\stopitemize

\stoptitle

\starttitle[title=Some things remain or adapt]

\startitemize
    \startitem Horizontal windmills for electricity become obsolete for vertical ones. \stopitem
    \startitem Tomorrows batteries will make today's batteries look ridiculous. \stopitem
    \startitem The last few years \TEX\ hasn't changed that much and has demonstrated adaptation. \stopitem
    \startitem Usage of \TEX\ behind the scenes is definitely usage. \stopitem
    \startitem This also depends on long term commitment, dedication to quality, and a wish for control. \stopitem
    \startitem (And let's not forget the fun part.) \stopitem
\stopitemize

\stoptitle

\starttitle[title=Some observations]

\startitemize
    \startitem In the meantime \TEX\ and installations are quite stable and usable. \stopitem
    \startitem What was once large is now small. System demands are rather low. \stopitem
    \startitem There are plenty of examples around to get one going. \stopitem
    \startitem There are enough alternatives to keep those who dislike it happy. \stopitem
    \startitem We don't need to catch up with every temporary web driven \quote {innovation}. \stopitem
    \startitem The users of \CONTEXT\ are dedicated to quality and flexibility. \stopitem
    \startitem Instead \quote {making \TEX\ great again} we can focus on \quote {keeping \TEX\ little}. \stopitem
\stopitemize

\stoptitle

\starttitle[title=Example]

If we have time: The hidden change in fundamentals from \MKII\ to \MKIV\ to
\MKXL: parameters, performance, etc. Did anyone (except the few involved in
dealing with it) notice it?

\stoptitle

\starttitle[title=The roadmap]

\startitemize
    \startitem We now have the extended \TEX\ engine that we always wanted. \stopitem
    \startitem There is not much that needs to be added and it's opened up anyway. \stopitem
    \startitem The (now) third major \CONTEXT\ upgrade is probably as good as we can make it. \stopitem
    \startitem We can and will play with the builders in the engine to see where we can improve and add. \stopitem
    \startitem We can however focus on the tricky bits of automated rendering. \stopitem
    \startitem We can try to squeeze out some performance but we also want to keep the design. \stopitem
\stopitemize

But:

\startitemize
    \startitem It are the users who have to come up with demands. We anyway have some. \stopitem
\stopitemize

\stoptitle

\stopdocument

